

%\newcommand*{\ACM}{}%

\ifdefined\ACM

%\documentclass[sigplan,screen]{acmart}
\documentclass[manuscript,screen,review]{acmart}

\else
 \documentclass[12pt]{extarticle}
%\usepackage{libertine}
\usepackage{cuted}%\usepackage{widetext}
../synthesis/usepackage.tex
\usepackage{cancel}
\usepackage{tikz-3dplot}
\usepackage{subcaption}
\addbibresource{./sample.bib}

\fi
\usepackage{physics} 

\begin{document}
../synthesis/newcommands.tex
\title{ A Note on an Old Proof of 'Quantum Toom's Rule' }
\author{David Ponarovsky}
\maketitle

% Constant depth fault tolerance construction.   
\newcommand*{\xb}{\mathbf{x}}

\ifdefined\SWEEPHEA
\else 
  \newcommand*{\SWEEPHEA}{}
  %\newcommand{\History}[1]{\mathbf{History}(#1)}
\newcommand{\History}[1]{H_{#1}}
  %\newcommand{\Time}[1]{\mathbf{Time}(#1)}
  \newcommand{\Time}[1]{T(#1)}
  \newcommand{\Excuse}[1]{\mathbf{Excuse}(#1)}
  \newcommand{\Excusei}[1]{\mathbf{Excuse}_{i}(#1)}
  \newcommand{\Size}[1]{\mathbf{Size}(#1)}
  \newcommand{\Span}[1]{\mathbf{Span}(#1)}
  \newcommand{\pre}[1]{\text{pre}(#1)}
  \newcommand{\poli}[1]{\text{pole}_{i}(#1)}
  \newcommand{\prei}[1]{\text{pre}_{i}(#1)}
  \newcommand{\predd}[1]{\text{pre}_{d+1}(#1)}
  \newcommand{\rotf}[2]{\theta_{#1 ,#2}}
  \newcommand*{\rotfs}{\rotf{v}{S^{\prime}}}
  \newcommand{\rotfv}[2]{\mathbf{\rotf{#1}{#2} }}
  \newcommand*{\rotfvs}{\mathbf{\rotfs{} }}
\fi



\begin{abstract}
This note examines the correctness of a proof of the 'Quantum Toom's rule' given in the old thesis \cite{ahn2004extending}, written almost twenty-five years ago. Although the presented ideas seem to point in the right direction, we found gaps that undermine the solidity of the work. Here, we present three of them. In a separate accompanying paper, we provide a candidate for a correct proof of the main theorem.
\end{abstract}
\section{Introduction} 
The quantum Toom's rule has been mentioned extensively in the literature \cite{Dennis_2002}, yet the question of the existence of a rigorously correct proof is not a consensus. One promising candidate is the work done by \cite{ahn2004extending} in her Ph.D. thesis. We emphasize that her work is meaningful and contains clever concepts regardless of whether it's correct as a whole. Yet, as recent works have been relying on her result, it becomes obligatory to understand if the 'Quantum Toom's rule', as stated by her, is indeed valid. Here, we present three counterexamples to three arguments raised in her proof. The goal of the note is to be short and aimed, and therefore the text assumes that the reader is fluent in the proof details.

\section{Failure of the Disobedient-to-Obedient Construction}

We first examine the conversion from a disobedient to an obedient explanation
proposed in \cite[pp.~105--108]{ahn2004extending}. The construction is supposed to
replace an actual error history by one satisfying the prescribed oriented
predecessor relation, while introducing new noise sources only at time steps
that already contain physical noise. The following example exhibits an
obstruction to the stated sweep-and-erase procedure.

Work in the four-dimensional cubic lattice, and restrict attention to the
hyperplane $w=0$. Let $P_{ij}(x,y,z)$ denote the elementary $ij$-plaquette
with lower corner $(x,y,z,0)$, where $i,j\in\{x,y,z\}$. All chains below
are over $\mathbb F_2$, so addition means symmetric difference. Introduce
three faults at time $0$:
\begin{equation}\label{eq:three-faults}
 a=P_{xy}(0,0,1),\qquad
 b=P_{xz}(0,1,0),\qquad
 c=P_{xz}(0,1,1),\qquad E_0=a+b+c.
\end{equation}
There are no further physical faults. Define
\begin{equation}\label{eq:two-produced-plaquettes}
 p=P_{yz}(1,0,0),\qquad q=P_{yz}(0,0,0).
\end{equation}
Let $U(E)$ be the set of plaquettes flipped by one synchronous, noiseless
application of Ahn's rule: a plaquette is flipped if its two positive edges
belong to $\partial E$. A direct evaluation of the syndrome gives
\begin{equation}\label{eq:three-fault-evolution}
 U(E_0)=a+b+c+p+q,\qquad E_1=E_0+U(E_0)=p+q,
\end{equation}
followed by
\begin{equation}
 U(E_1)=p+q,\qquad E_2=0.
\end{equation}
Thus a perfectly valid, noiseless decoding step creates the erroneous
plaquette $p$, although $p$ was not erroneous previously.

Consider the prescribed set of \emph{obedient predecessors} of $p$: these
are $p$ itself and the neighbouring plaquettes that share a positive edge
of $p$ as one of their negative edges. The positive $z$-edge of $p$
(the edge parallel to $y$ at $z=1$) carries syndrome because of $a$;
the positive $y$-edge of $p$ (parallel to $z$ at $y=1$) carries syndrome
because of $b$. However, the first shared edge is a \emph{positive}
$x$-edge of $a$, and the second is a \emph{positive} $x$-edge of $b$.
Consequently neither $a$ nor $b$ is an obedient predecessor of $p$.
Neither is $c$. Indeed, every permitted predecessor other than $p$
has center $x$-coordinate at least $1$, whereas $a,b,c$ have center
$x$-coordinate $1/2$. In particular, if $\mathcal P(p)$ denotes the
prescribed predecessor neighbourhood,
\begin{equation}\label{eq:no-obedient-predecessors}
 \mathcal P(p)\cap E_0=\varnothing.
\end{equation}

Let $T=\{a,b,c,p,q\}$. These are all the plaquettes that occur either
in the initial error or among the flips in its entire noiseless
continuation. The stated conversion draws previous explanation cells and
Toom-flipped cells, and subsequently \emph{erases} cells; its operations
do not introduce any additional plaquette outside $T$. But
\begin{equation}
 \mathcal P(p)\cap T=\{p\},
\end{equation}
so at most one of the required two distinct predecessors is available.
Moreover, $p$ cannot simply be discarded if the new explanation is
required to reproduce the syndrome at time $1$: the edge
\begin{equation}
 e=\{1\}\times[0,1]\times\{0\}\times\{0\}
\end{equation}
belongs to $\partial E_1$, and $p$ is the only plaquette in $T$
incident to $e$. Hence any subset of $T$ with boundary $\partial E_1$
necessarily contains $p$. Declaring $p$ a new noise source at time $1$
would instead introduce an obedient noise source at a time step with
no physical noise. Thus the stated operations cannot simultaneously
preserve the time-$1$ syndrome, provide two obedient predecessors of
$p$, and preserve the claimed noise-time correspondence.

\paragraph{Amplification.}
Place $m$ translated copies of \eqref{eq:three-faults} at mutual
distance at least $5$ (on a sufficiently large torus, or in the
infinite lattice), all at the same initial time. The copies do not
interact under the two noiseless updates above. There are exactly
$3m$ physical faults at time $0$, no later faults, and $m$ distinct
plaquettes $p_1,\ldots,p_m$ appearing at time $1$. For each $p_j$,
the corresponding drawn set contains at most one permitted predecessor,
while an edge of $\partial E_1$ forces $p_j$ to remain in any
syndrome-preserving representative constructed from those cells.
Thus the same obstruction occurs independently at arbitrarily many
locations. This refutes the sufficiency of the \emph{stated conversion
procedure}; it does not rule out a different conversion that adds
auxiliary plaquettes and proves a new charging bound relating its
noise sources to actual faults.

\section{Dense-Collection Probability Bound Relies on an Unjustified Completion Argument}

The deduction of Eq.~(7.20) in \cite[pp.~109--110]{ahn2004extending} uses the
principle that, if every configuration in an event $A$ can be
completed into an event $B$ by adding at most $k$ faults, then
$\Pr(A)\le\varepsilon^{-k}\Pr(B)$. Such a conclusion does not follow
from the existence of completions alone, even for independent noise
of a fixed arbitrarily small rate. We give a rectangular example.

Let $R$ be a rectangle containing $N$ fault locations, where
$4\mid N$; for instance, take a $4m\times m$ rectangle, so
$N=4m^2$. Each location is faulty independently with probability
$\varepsilon$, where $\varepsilon=1/100$ is fixed, independently of $N$.
Define the events
\begin{align}
 A&=\{\text{at least $N/4$ locations of $R$ are faulty}\},\\
 B&=\{\text{all $N$ locations of $R$ are faulty}\}.
\end{align}
Every configuration in $A$ can be completed to one in $B$ by
introducing at most $3N/4$ additional faults. Nevertheless,
\begin{align}
 \Pr(A)
 &\ge {N\choose N/4}\varepsilon^{N/4}
             (1-\varepsilon)^{3N/4},\\
 \Pr(B)&=\varepsilon^N.
\end{align}
The alleged probability division would imply
$\Pr(A)\le\varepsilon^{-3N/4}\Pr(B)=\varepsilon^{N/4}$.
In fact,
\begin{equation}\label{eq:completion-multiplicity}
 \frac{\Pr(A)}{\varepsilon^{-3N/4}\Pr(B)}
 \ge {N\choose N/4}(1-\varepsilon)^{3N/4}
 \ge \bigl[4(1-\varepsilon)^3\bigr]^{N/4}>1.
\end{equation}
Here we used ${4r\choose r}\ge 4^r$ for $r=N/4$, and the
last inequality holds at the fixed constant $\varepsilon=1/100$.
The discrepancy grows exponentially with the area of $R$.

The obstruction is multiplicity: many partially faulty configurations
can have the same completed configuration. A valid probability
comparison requires a specified map from original to completed
noise histories, together with a bound on its number of preimages
(or an equivalent weighted-counting estimate). This rectangular
example disproves the general probability-division \emph{inference},
not Eq.~(7.20) itself. In Ahn's geometric setting the claimed
completion cost is $3L/4$, where $L$ is a \emph{side length}, not
$3N/4$ for a rectangle of area $N$. Establishing that particular
bound additionally requires the appropriate geometric completion
and its multiplicity estimate.

\section{Toom's Rule in Higher Dimensions Does Not Shrink All Droplets}

In four dimensions Ahn flips a plaquette if both of its positive
edges carry syndrome \cite[p.~105]{ahn2004extending}. For qubits on $k$-dimensional
cells, the direct all-positive-facets extension is to flip a $k$-cell
if all its $k$ positive $(k-1)$-facets belong to the syndrome.
We show that this natural rule has a contractible, nonzero fixed
point for every $k\ge3$. In particular, it fails for qubits on
three-faces in six dimensions. The thesis does not give an equally
explicit higher-dimensional update prescription; accordingly, the
example applies to this natural extension rather than to every
possible higher-dimensional decoding rule.

In the six-dimensional cubic lattice with coordinates
$(x,y,z,w,u,v)$, put qubits on three-cells and syndrome checks on
square faces. Let
\begin{align}
 A&=[0,1]_x\times[0,1]_y\times[-1,0]_z,
   &&w=u=v=0,\\
 B&=[0,1]_x\times[0,1]_y\times[-1,0]_w,
   &&z=u=v=0.
\end{align}
Their intersection is the square
\begin{equation}
 Q=[0,1]_x\times[0,1]_y,
 \qquad z=w=u=v=0,
\end{equation}
which is a positive facet of \emph{both} cubes. For the error chain
$E=A+B$, its syndrome $\sigma=\partial E$ does not contain $Q$,
since the two incidences cancel over $\mathbb F_2$.

Each of $A$ and $B$ therefore has only two of its three positive
facets in $\sigma$, so neither is flipped. Nor is any other
three-cell $C$: two distinct elementary cubes share at most one
square facet; hence $C\notin\{A,B\}$ can have at most one positive
syndrome facet contributed by $A$ and at most one contributed by
$B$. In particular, it has at most two positive facets in
$\sigma$, whereas a flip requires three. Thus the correction
chain is empty and
\begin{equation}\label{eq:fixed-six-dimensional-droplet}
 U(\partial E)=0,\qquad E_t=E\quad\text{for every }t\ge0.
\end{equation}
The union $A\cup B$ is contractible, and $\partial E\ne0$. Hence
this is a finite, contractible error droplet with nonzero syndrome
that never shrinks under noiseless evolution.

More generally, take two elementary $k$-cubes sharing a common
$(k-1)$-face that is positive for both. Their common facet cancels
from the syndrome, leaving only $k-1$ positive syndrome facets on
each erroneous cube. Any other $k$-cube meets each of the two
original cubes in at most one $(k-1)$-facet, so it has at most
two positive syndrome facets. For $k\ge3$, no cell is flipped.
The same fixed point therefore exists in every ambient dimension
$d\ge k+1$, in particular in the $d=2k\ge6$ family.

%\inb{\textbf{TODO:}}
%\begin{enumerate}
%
%
%  \item Enter that we care only about pauli errors somewhere. 
%  \item Rewrite the all outline. 
%
% % \item Define what is mean a finite chain. 
%
%  \item Mention that in 'applying the sweep rule on a chain', we refer to the application of the first stage of the sweep cycle. Namely, all the vertices work, yet we don't care about the dual base.
%  \item In all the claims about chains, set $x$ to note the default chain (instead $z$ which we keep to stabilizers) 
%  \item Change all the 'Pauli' to capital letter. 
%  \item Add after the definition of the computation DAG, that the order $J$ has to respect time order.  
%  %\item Enter the \textbf{Str} notation to the agreement claim. \Cref{claim:agreement}. 
%  \item Using \Cref{claim:polesmove} in \Cref{claim:shrinkslice} needs some work in the case that the slice contains multiple connected components. 
%\end{enumerate}



\printbibliography

%\section{Garbage}


Now we are ready to prove that any error is equivalent, up to addition by stabilizer, to an error whose maximal $L_{d+1}$ vertex (or its edge) sees a syndrome (and that maximum is equal to or lower than the maximal value of $L_{d+1}$ in the original error). \Cref{fig:edge} visualize the idea. \Cref{claim:reduce_stab} introduces a 'single application' of \Cref{claim:agreement} to 'advance' the corner (the vertex that gets the maximal value of $L_{d+1}$). The corner might not be unique, so in the proof, we advance all the corners simultaneously. Yet, it might occur that the yielded assignment still doesn't admit corners that see a syndrome. Thus, \Cref{claim:extermum} repeats  'corners advancing' steps on until it ensures the desired outcome - that the corner vertex of the final yielded assignment sees a syndrome. 
\begin{claim}
  \label{claim:reduce_stab}
  Let $z$ be a finite assignment. Suppose that 
  \begin{equation*}
    \begin{split}
            \max \big\{ L_{d+1}(u) :  u \in z \big\} > \max \big\{ L_{d+1}(u) : u \in \partial^{-} z \big\}
    \end{split}
  \end{equation*}
  Then there is $\tilde{z} \sim_{\mathcal{C}_{X}} z$ such that:
\begin{equation*}
  \begin{split}
    \max \left\{ L_{d+1}(u) : u \in \tilde{z} \right\} & \le  \max \big\{ L_{d+1}(u) : u \in z \big\} - 1
  \end{split}
\end{equation*}
\end{claim}
\begin{proof}
  Let $v_{1}, v_{2}, \ldots, v_{k}$ be the vertices in $z$ that achieve the maximum for $L_{d+1}$, namely $v_{1}, v_{2}, \ldots, v_{k} = \arg \max_{v \in z} L_{d+1}(v)$, and assume that $v_{i} \notin \partial^{-} z$ for any $i \in [k]$. (Notice that the maximum is indeed achieved since $z$ is finite.)

Consider the $i$th vertex $v_{i}$. By definition, since $v_{i} \notin \partial^{-} z$, we have $\partial^{-} z|_{v_{i}} = 0$, and since $v_{i}$ is a maximum point of $L_{d+1}$ in $z$, it is also a local maximum, and therefore $z|_{v_{i}+} = z|_{v_{i}}$. 
Combining the two, we get that $z|_{v_{i}}$ is a codeword of the small code at $v_{i}$, namely $z|_{v_{i}} \in \mathcal{C}_{v_{i}}$. Thus, by \Cref{claim:agreement}, there is a stabilizer $w_{i} \in \text{Stab}_{v_{i}}$ such that $w_{i}$ and $z$ agree on the positive faces of $v_{i}$, namely $w_{i}|_{v_{i}} = z|_{v_{i}}$. 

Set $\tilde{z} \leftarrow z + \sum_{i}w_{i}$. Since $\sum_{i}w_{i}$ is a stabilizer, we have $\tilde{z} \sim_{C_{X}} z$. In addition, for any $j \neq i$, it follows that $v_{j} \notin w_{i}$.
Otherwise, there exists $S^{\prime} \subset S$ such that $v_{j} \in \rotf{v_{i}}{S^{\prime}}/ v_{i}$, and therefore $L_{d+1}(v_{j}) > L_{d+1}(v_{i})$, in contradiction to $v_{i}$ and $v_{j}$ being both maximum points for $L_{d+1}$ in $\partial^{-} z$. 
Thus, we have:
\begin{equation*}
  \begin{split}
    \tilde{z}|_{v_{i}} = (z + \sum_{j}w_{j})|_{v_{i}} =  (z + w_{i})|_{v_{i}} = 0
  \end{split}
\end{equation*}
Namely, for any $i$, the $i$th vertex $v_{i}$ is not in $\tilde{z}$. Hence, the vertex domain of $\tilde{z}$ is contained in the union of the vertices of $z$ with the vertices of $w_{1}, w_{2}, \ldots, w_{k}$, substituting $v_{1}, v_{2}, \ldots, v_{k}$. Overall, we get:
\begin{equation*}
  \begin{split}
  \max \left\{ L_{d+1}(u) : u \in \tilde{z} \right\} & \le \max \big\{ L_{d+1}(u) : u \in z \bigcup_{i} w_{i} \big/ \bigcup_{i} v_{i} \big\}  \\ 
     & \le  \max \big\{ L_{d+1}(u) : u \in z \big\} - 1
  \end{split}
\end{equation*}
\end{proof}



\begin{claim}
  \label{claim:extermum}
  Let $z$ be a finite assignment. Then there is $\tilde{z} \sim_{\mathcal{C}_{X}} z$ such that 
  \begin{equation*}
    \begin{split}
      \max \big\{ L_{d+1}(u) :  u \in \tilde{z} \big\} = \max \big\{ L_{d+1}(u) : u \in \partial^{-} z \big\}
    \end{split}
  \end{equation*}
\end{claim}
\begin{proof}
Let $z$ be a finite assignment, and let $k$ be an integer such that $k \ge 1$ and
  \begin{equation*}
    \begin{split}
            \max \big\{ L_{d+1}(u) :  u \in z \big\} =  \max \big\{ L_{d+1}(u) : u \in \partial^{-} z \big\} + k
    \end{split}
  \end{equation*}
  For an integer $l \le k$, consider the assignments $\tilde{z}_{0},\tilde{z}_{1},..,\tilde{z}_{l}$ defined as follow: 
  \begin{itemize}
    \item The first element, $\tilde{z}_{0} = z$. 
    \item If the $i-1$th element satisfies 
  \begin{equation*}
    \begin{split}
      (\star ) \ \ \max \big\{ L_{d+1}(u) :  u \in \tilde{z}_{i-1} \big\} > \max \big\{ L_{d+1}(u) : u \in \partial^{-} \tilde{z}_{i-1} \big\} 
    \end{split}
  \end{equation*}
      Then define the $i$th element, $\tilde{z}_{i}$, to be the result of applying the construction, given in \Cref{claim:reduce_stab}, on $\tilde{z}_{i-1}$ , namely  $\partial \tilde{z}_{i+1} = \partial z_{i}$ and  
\begin{equation*}
  \begin{split}
    \max \big\{ L_{d+1}(u) :  u \in \tilde{z}_{i} \big\} \le  \max \big\{ L_{d+1}(u) : u \in \tilde{z}_{i-1} \big\} - 1
  \end{split}
\end{equation*}
If the $i-1$th element doesn't satisfies $\star$, then set $l = i -1$.
\item If the $k$th element was defined, set $l = k$.
  \end{itemize}
Since $\sim_{C_{X}}$ is a transitive relation, we have that for any $i \in [l]$, $z \sim_{C_{X}} \tilde{z}_{i}$. If $l \neq k$, then there exists $\tilde{z}_{i}$ such that $\tilde{z}_{i} \sim_{C_{X}}$ and $\tilde{z}_{i}$ satisfies $\star$, namely by setting $\tilde{z} = \tilde{z}_{i}$, we get the desired.
So, it is left to handle the case where $l = k$. In that case, set $\tilde{z} = \tilde{z}_{k}$. By \Cref{claim:reduce_stab}:
\begin{equation*}
  \begin{split}
    \max \big\{ L_{d+1}(u) :  u \in \tilde{z} \big\} \le  \max \big\{ L_{d+1}(u) : u \in \partial^{-} z \big\}
  \end{split}
\end{equation*}
Yet, since $\tilde{z} \sim_{C_{X}} z$, it follows that $\partial \tilde{z} = \partial z$, So: 
\begin{equation*}
  \begin{split}
     \max \big\{ L_{d+1}(u) :  u \in \partial^{-} z \big\} = \max \big\{ L_{d+1}(u) :  u \in  \partial^{-} \tilde{z} \big\} \le \max \big\{ L_{d+1}(u) : u \in z \big\} 
  \end{split}
\end{equation*}
Therefore, we have the equality:
\begin{equation*}
  \begin{split}
    \max \big\{ L_{d+1}(u) :  u \in \tilde{z} \big\} = \max \big\{ L_{d+1}(u) : u \in \partial^{-} z \big\}
  \end{split}
\end{equation*}
\end{proof}


 
\end{document}
